\subsection{Sensitivity to Independent PMU Noise}
\label{sec:noise-stress}

To examine sensitivity beyond the calibrated noise, we reran the
nominal bank with the same equations and publication thresholds on
four held-out calibration event records (\RevNoiseLoadEvents{} load and
\RevNoiseLineEvents{} line interventions, matched network) and two 600-s
normal records (\RevNoiseNormalHours\,h). Six noise scenarios were
generated with the simulator's own noise stream: the original realization
scaled by $\times1$, $\times5$, $\times10$ and $\times20$, and the original
plus independent per-PMU relative-magnitude and angle noise, each with
standard deviation $\sigma=0.0005$ or $0.001$ (angles in radians).
The added perturbation has approximately \RevNoiseAddedRmsLow{}\% or \RevNoiseAddedRmsHigh{}\% RMS total
vector error (TVE), respectively, before combination with the original
noise. Voltage-angle noise is propagated to frequency and ROCOF by the
simulator's 21-report estimator. The $\times1$ scenario reproduces the
original predictions message for message. The bank re-estimates its noise
scale and observation discrepancy from the noisy 300-report prefix; the
operating point is kept at its nominal-noise estimate because derivative
matrices exist only there. The implementation refuses to start when the
estimated noise scale leaves its calibrated support of $[0.1,10]$; for this
test only, that bound was widened, and every other line of code is
unchanged.

\RevNoiseDiscussion{}
